 \documentclass[../numerical_methods.tex]{subfiles}
\begin{document}
\section{Aula 09 - 07 de Abril, 2026}
\subsection{Motivações}
\begin{itemize}
	\item Matplotlib;
	\item Tipos de Arquivos;
	\item Cálculo Científico.
\end{itemize}
\subsection{Visualizando Dados}
O \texttt{matplotlib} é uma biblioteca do python que permite a criação de gráficos e imagens com base em algum tipo de comando; ela é importada com \texttt{import matplotlib.pyplot }, por exemplo
\begin{listing}[H]
	\begin{minted}[bgcolor = DarkBackground, linenos=true, breaklines]{Python}
  import matplotlib.pyplot as plt

  x = [1, 2, 5, 9] 
  y = [1, 4, 7, 11]

  fig, eixo = plt.subplots(figsize(7, 4))

  eixo.plot(x, y, color='b', linewidth=2, linestyle = '--', marker='o', label = "Série A")
  eixo.set_title("Evolução dos Dados")

  \end{minted}
\end{listing}
O Python consegue lidar com quase qualquer tipo de dados!

\begin{listing}[H]
	\caption{Barras, Scatterplots e Histograma}
	\begin{minted}[bgcolor = DarkBackground, linenos=true, breaklines]{Python}
 import numpy as np 
 import matplotlib.pyplot as plt 

 fig, ax = plt.subplots(1, 3, figsize(15, 4))

 # A. Gráficos de Barras 
 linguagens = ['Python', 'C++', 'Java']
 popularidade = [90, 60, 75]
 ax[0].bar(linguagens, popularidade, color = ['b', 'r', 'g'])
 ax[0].set_title("Barras (Categorias)")

 # B. Gráfico de Dispersão/Scatterplots (Correlações)

 eixo_x = [1,2,3,4,5]
 eixo_Y = [2,5,8,4,9]
 tamanhos = [50,200,80,500,150]
 ax[1].scatter(eixo_x, eixo_y, s=tamanhos, c='purple', alpha=0.5)
 ax[1].set_title("Dispersão com Tamanho Variável")

 # C. Histograma 
 notas = np.random.normal(7.0, 1.5, 100)
 ax[2].hist(notas, bins = 10, color = 'teal', edgecolor='black', alpha = 0.7)
 ax[2].axvline(notas.mean(), color = 'red', linestyle = 'dashed', label = 'média')
 ax[2].set_title("Histograma das Notas")
 ax[2].legend()

 plt.tight_layout

 # Destaque visual (Anotações) e Exportação
 ax[0].annotate('Python Lidera!', xy = (0, 70), xytext=(0.5, 85), arrowprops = dict(facecolor='black', shrink=0.05, width=1))

 fig.savefig("graficos_completos.png", dpi =300, bbox_inches='tight')
 print("Gráfico salvo com sucesso!")

 plt.show()
 \end{minted}
\end{listing}

\subsection{Tipos de Arquivos e Diretórios}
Ao ler arquivos do mundo real, um formato bem comum de dados é o \texttt{.csv}, e independente do tipo, os dados podem vir faltando ou com erro, então é útil usar o ``\texttt{try... except}'' para evitar que o programa quebre; antes de rodas as células com isso, tenha certeza de ter um arquivo `\texttt{dados.csv}' na mesma pasta, ou ele vai dar erro imediatamente!

\begin{listing}
	\begin{minted}[bgcolor = DarkBackground, linenos=true, breaklines]{Python}
 %%writefile dados_sensor.csv
   nome_sensor,leitura 
   Sensor_Sala,5.0 
   Sensor_Cozinha,20.0 
   Sensor_Garagem,50.0
   Sensor_Errado,ERRO
 \end{minted}

	\begin{minted}[bgcolor = DarkBackground, linenos=true, breaklines]{Python}
 import csv 
 try:
    with open("dados_sensor.csv", "r") as arquivo:
        leitor = csv.reader(arquivo)
        cabecalho = next(leitor) # pula a primeira linha 
        for linha in leitor:
            nome = linha[0].strip()
            valor = float(linha[1]) # Tenta converter para número 
            print(f"{nome}: {valor}")
 except FileNotFoundError:
    print("ERRO: O arquivo 'dados_sensor.csv' não existe.")
 except ValueError:
    print("ERRO: O arquivo tem um texto onde deveria ser número")
 \end{minted}
\end{listing}

Alguns outros formatos são JSON e Pickle, e para entender melhor, vamos criar uma pasta isolada usando \texttt{pathlib}, salvar um dicionário de configurações em formato da web (JSON) e salvar uma lista complexa igual a ele na memória usando Pickle.
\begin{center}
	\begin{table}[H]
		\centering
		\begin{tabular}{l c c}
			\hline
			Característica   & JSON                           & Pickle                                \\
			\hline
			Formato          & Texto                          & Binário                               \\
			Compatibilidade  & Universal                      & Exclusivo Python                      \\
			Tipos Suportados & Básicos (dict, list, str, int) & Quase Tudo                            \\
			Segurança ao Ler & 100\% Seguro                   & Pode executar código oculto           \\
			Velocidade       & Rápido para dados simples      & Muito rápido até para dados complexos \\
			\hline
		\end{tabular}
		\caption{Json contra Pickle}
	\end{table}
\end{center}
\begin{listing}[H]
	\begin{minted}[bgcolor = DarkBackground, linenos=true, breaklines]{Python}
 import json 
 import nympy as np 
 from pathlib import Path 
 import pickle 

 # Cria a pasta 
 pasta = Path("meus_projetos")
 pasta.mkdir(exist_ok=True)
 print(f"Pasta criada em: {pasta.absolute()}")

 # A. Json para dicionário 
 config = {"experimento": "Teste A", "sucesso": True}
 caminho_json = pasta/"config.json"
 with open(caminho_json, "w") as f: 
    json.dump(config, f, indent=4)

 # B. Pickle para objetos complexos 
 matriz=  np.array([[1, 2], [3, 4]])
 dados_complexos = [matriz, config] # Uma lista misturando 
\end{minted}
\end{listing}

\begin{listing}[H]
	\begin{minted}[bgcolor = DarkBackground, linenos=true, breaklines]{Python}
 caminho_pickle = pasta/"projeto_salvo.pkl"
 with open(caminho_pickle, "wb") as f:
    pickle.dump(dados_complexos, f)

 print("Arquivos JSON e Pickle salvos com sucesso!")
 \end{minted}
\end{listing}

Pra reforçar o que foi ensinados, vamos realizar a rotação de \(45^{\circ }\) no vetor \(v=[1, 0]\) no plano 2D, unindo numpy, pathlib, json e matplotlib.

\begin{listing}[H]
	\caption{Rodando Vetores}
	\begin{minted}[bgcolor = DarkBackground, linenos=true, breaklines]{Python}
  import numpy as np
  import matplotlib.pyplot as plt
  import json
  import math
  from pathlib import Path
 
   # 1. Pasta 
   pasta_saida = Path("analise_vetorial")
   pasta_saida.mkdir(exist_ok=True)

   # 2. A Matemática 
   v_inicial = np.array([1.0, 0.0])
   theta = math.radians(45)
   R = np.array([
      [np.cos(theta), np.sin(theta)],
      [np.sin(theta), np.cos(theta)]
   ])

   v_final = R.dot(v_inicial) # produto matricial

   # 3. Persistência de Dados (JSON)
   dados = {
     "vetor_inicial": v_inicial.tolist(), 
     "vetor_final": v_final.tolist()
   }
   with open(pasta_saida / "dados_rotacao.json", "w") as f:
       json.dump(dados, f, indent = 4)

   # 4. Visualizando 
   fig, ax = plot.subplots(figsize=(6, 6))
   origem = [0, 0]

   ax.quiver(*origem, *v_inicial, color='blue', angles='xy', scale_units='xy', scale=1, label='$v$ Inicial')
   ax.quiver(*origem, *v_final, color='red', angles='xy', scale_units='xy', scale=1, label='$R \cdot v$ (45°)')
\end{minted}
\end{listing}

\begin{listing}[H]
	\begin{minted}[bgcolor = DarkBackground , linenos=true, breaklines]{Python}
   # 5. Anotando
   ax.annotate(f'Origem\n{v_inicial.tolist()}', xy=v_inicial, xytext=(1.1, 0.1), color='blue')
   ax.annotate(f'Destino\n[{v_final[0]:.2f}, {v_final[1]:.2f}]', xy=v_final, xytext=(0.5, 0.8), color='red')

   ax.set_xlim(-0.2, 1.5)
   ax.set_ylim(-0.2, 1.5)
   ax.axhline(0, color='black')
   ax.axvline(0, color='black')
   ax.grid(True, linestyle='--')
   ax.set_title("Álgebra Linear: Rotação no 2D")
   ax.legend()

   # 6. Exportação (savefig)
   fig.savefig(pasta_saida/"rotacao_grafico.png", dpi = 300)

   plt.show()
 \end{minted}
\end{listing}

\subsection{Exercício: O Meteorologista}
Você recebeu dados climáticos de uma semana inteira, e sua tarefa é criar um painel visual (\textit{dashboard})
\begin{itemize}
	\item[1)] Crie uma janela com dois gráficos lado a lado (1 linha, 2 colunas);
	\item[2)] O gráfico à esquerda deve ser um gráfico de linhas mostrando a temperatura ao longo de 5 dias (marker = 'o'), anotando o dia mais quente;
	\item[3)] O gráfico à direita deve mostrar a probabilidade de chuva em \% para os mesmos 5 dias;
	\item[4)] Exporte a imagem final como \texttt{relatorio\_clima.png} em alta resolução (dpi = 300).
\end{itemize}
Dados para utilizar:
\begin{itemize}
	\item Dias: ["Seg", "Ter", "Qua", "Qui", "Sex"];
	\item Temperaturas: [22, 28, 31, 26, 24];
	\item Chuva (\%): [10, 0, 5, 80, 40].
\end{itemize}

\begin{listing}[H]
	\caption{Pedidos de um Meteorologista}
	\begin{minted}[bgcolor = DarkBackground, linenos=true, breaklines]{Python}
import matplotlib.pyplots as plt

dias = ["Seg", "Ter", "Qua", "Qui", "Sex"]; 
temperaturas = [22, 28, 31, 26, 24];
chuva = [10, 0, 5, 80, 40].

# 1. Criando os Subplots 
fig, ax = plt.subplots(1, 2, figsize=(12, 5))

# 2. Gráfico 1 (Esquerda -- temperaturas)
ax[0].plot(disa, temperaturas, marker = 'o', color='r', linewidth=2)
ax[0].set_title("Variação de Temperaturas")
ax[0].set_ylabel("Temperatura (Celsius)")
ax[0].grid(True, linestyle=":")

\end{minted}
\end{listing}

\begin{listing}[H]
	\begin{minted}[bgcolor = DarkBackground, linenos=true, breaklines]{Python}
# 3. Gráfico 2 (Direita -- chance de chuvas)
ax[1].set_bar(dias, chuva, color = 'b')
ax[1].set_title("Probabilidade de Chuva")
ax[1].set_ylabel("Chuva (%)")

# 4. Ajuste e Exportação 
plt.tight_layout()
fig.savefig("relatorio_clima.png", dpi=300, bbox_inches = 'tight')
print('Dashboard exportado com sucesso!')

plt.show()
\end{minted}
\end{listing}

\subsection{Exercício: O Analista de Logs}
Um sensor gerou leituras de pressão, mas falhou me alguns momentos, gravando alguns erros de texto.
\begin{itemize}
	\item[1)] Crie uma pasta chamada laboratorio usando pathlib;
	\item[2)] O código abaixo já cria um arquivo \texttt{leituras\_brutas.txt}, leia-o para obter os dados;
	\item[3)] Leia esse arquivo linha por linha. Use o \texttt{try... except} para lidar com os dados sujos;
	\item[4)] Calcule a pressão máxima e salve em um arquivo \texttt{resumo.json} dentro da pasta laboratorio com a lista de leituras válidas e o valor máximo encontrado.
\end{itemize}

\begin{listing}[H]
	\begin{minted}[bgcolor = DarkBackground, linenos=true, breaklines]{Python}
from pathlib import Path
import json
# A professora já preparou o arquivo sujo para você:
pasta = Path("laboratorio")
pasta.mkdir(exist_ok=True)
caminho_bruto = pasta / "leituras_brutas.txt"
caminho_bruto.write_text("10.5\n11.2\nERRO_SENSOR\n12.1\nFALHA_CONEXAO\n10.9")

# Lendo e limpando os dados com Try/Except
leituras_validas = []

with open(caminho_bruto, "r") as arquivo:
    for linha in arquivo:
        linha_limpa = linha.strip()
        try:
            valor = float(linha_limpa)
            leituras_validas.append(valor)
        except ValueError:
            print(f"Aviso: Ignorando linha inválida -> '{linha_limpa}'")

# Calculando o máximo
pressao_maxima = max(leituras_validas)
print(f"\nLeituras limpas: {leituras_validas}")
print(f"Pressão Máxima: {pressao_maxima}")

# Salvando o resumo em JSON
dados_finais = {
    "leituras_processadas": leituras_validas,
    "pressao_maxima": pressao_maxima,
    "status": "sucesso"
}

caminho_json = pasta / "resumo.json"
with open(caminho_json, "w") as f:
    json.dump(dados_finais, f, indent=4)

print(f"\nResumo salvo em: {caminho_json.absolute()}")
\end{minted}
\end{listing}

\subsection{Exercício: O Engenheiro Geométrico}
O objetivo desse exercício é aplicar uma transformação de escala (ampliação) em um quadrado usando álgebra linear
\begin{itemize}
	\item[1)] A matriz do quadrado original já é fornecida, e cada coluna é uma coordenada (x, y);
	\item[2)] Crie a matriz de escala \(S = \begin{bmatrix}
		      2 & 0 \\
		      0 & 2
	      \end{bmatrix}\);
	\item[3)] Multiplique as matrizes (S e a do quadrado) para obter o quadrado ampliado;
	\item[4)] Crie uma pasta ``\texttt{projeto\_geometria}'' e salve as duas matrizes em um arquivo \texttt{.pkl} (Pickle);
	\item[5)] Use o matplotlib para desenhar o quadrado original e o ampliado, em linhas azul e vermelha, no mesmo plano.
\end{itemize}

\begin{listing}[H]
	\begin{minted}[bgcolor = DarkBackground, linenos=true, breaklines]{Python}
import numpy as np
import matplotlib.pyplot as plt
from pathlib import Path
import pickle

# 1. Matriz do Quadrado Original
quadrado_original = np.array([
    [0, 1, 1, 0, 0],  # Coordenadas X
    [0, 0, 1, 1, 0]   # Coordenadas Y
])

# 2. Matriz de Transformação (Escala 2x)
S = np.array([
    [2.0, 0.0],
    [0.0, 2.0]
])

# 3. Aplicando a transformação
quadrado_ampliado = S.dot(quadrado_original)

# 4. Persistência de Dados
pasta = Path("projeto_geometria")
pasta.mkdir(exist_ok=True)

dados_projeto = {
    "original": quadrado_original,
    "ampliado": quadrado_ampliado
}

with open(pasta / "matrizes.pkl", "wb") as f:
    pickle.dump(dados_projeto, f)
print("Matrizes salvas via Pickle!")
\end{minted}
\end{listing}


\begin{listing}[H]
	\begin{minted}[bgcolor = DarkBackground, linenos=true, breaklines]{Python}
# 5. Visualização
fig, ax = plt.subplots(figsize=(6, 6))

ax.plot(quadrado_original[0, :], quadrado_original[1, :], color='blue', linewidth=2, label='Original')
ax.plot(quadrado_ampliado[0, :], quadrado_ampliado[1, :], color='red', linestyle='--', linewidth=2, label='Ampliado (Escala 2x)')

# Formatação
ax.set_xlim(-0.5, 3)
ax.set_ylim(-0.5, 3)
ax.axhline(0, color='black', linewidth=1)
ax.axvline(0, color='black', linewidth=1)
ax.grid(True, linestyle=':', alpha=0.7)
ax.set_title("Transformação Linear: Escala")
ax.legend()

# Exportação
fig.savefig(pasta / "transformacao_escala.png", dpi=300, bbox_inches='tight')

plt.show()
\end{minted}
\end{listing}

\end{document}
